% arXiv submission: SYN-01
% Title: The Architecture of Compositional Formal Meaning: A Categorical
%        Framework for Structural Unification of Grammar, Semantics,
%        and Computation
% Author: A. Padavano (Independent researcher)
% Primary category: cs.LO (Logic in Computer Science)
% Cross-list: cs.CL, cs.PL, math.CT
% MSC 2020: 18C50, 03B65, 68Q55, 03G30
% License: CC BY 4.0
% Date: 2026

\documentclass[12pt,a4paper]{article}

\usepackage[utf8]{inputenc}
\usepackage[T1]{fontenc}
\usepackage{amsmath,amssymb,amsthm}
\usepackage{hyperref}
\usepackage[sort&compress,numbers]{natbib}
\usepackage{booktabs}
\usepackage{graphicx}
\usepackage{alltt}
\usepackage{microtype}
\usepackage[margin=1in]{geometry}

\hypersetup{
  colorlinks=true,
  linkcolor=blue,
  citecolor=blue,
  urlcolor=blue
}

\newtheorem{definition}{Definition}
\newtheorem{theorem}{Theorem}
\newtheorem{conjecture}{Conjecture}
\newtheorem{proposition}{Proposition}

\title{The Architecture of Compositional Formal Meaning:\\
A Categorical Framework for Structural Unification\\
of Grammar, Semantics, and Computation}

\author{A.~Padavano\\
  \textit{Independent researcher, Studium Generale ORGANVM}}

\date{2026}

\begin{document}

\maketitle

\begin{abstract}
This paper synthesises two prior investigations: one identifying eight
structural bridges connecting three traditions of formal semantics (logical,
computational, and natural language), and another establishing a fourfold
correspondence among grammars, types, proofs, and categories across the
Chomsky hierarchy. We ask whether these results unify into a single
categorical architecture of compositional formal meaning. The central claim
is that meaning, in its compositional and type-driven core, is the functorial
passage from a syntactic category to a semantic category, exhibiting the same
categorical structure whether the syntax is linguistic, logical, or
computational. This is structural unification in Klein's Erlangen Programme
sense---identifying a common mathematical schema across independent
traditions---not scientific unification requiring novel empirical predictions.
We develop the architecture through seven components: meaning as functor via
fibred categories, compositionality as functoriality, the Lambek calculus as
syntactic category, Montague's rule-to-rule hypothesis as natural
transformation, the DisCoCat framework assessed with honest empirical
limitations, compact closed categories with attention to the
symmetric/non-symmetric distinction, and a topos-theoretic perspective with a
worked example of anaphora resolution via presheaves. The unification claim is
stated as a five-way correspondence among derivation, proof, programme,
morphism, and meaning-composition, supported by a commutative diagram. We
address the distributional challenge: DisCoCat is reframed as categorical
proof of concept, transformer success is confronted without dismissal, and
testable predictions from the categorical architecture are identified. We
conclude that the compositional typed core of all three semantic traditions
admits genuine structural unification under a single functorial architecture.
\end{abstract}

\noindent\textbf{Keywords:} category theory, formal semantics,
compositionality, functoriality, Curry-Howard-Lambek correspondence, Lambek
calculus, Montague grammar, DisCoCat, compact closed categories, topos theory,
structural unification, Erlangen Programme

\bigskip

\section{Introduction}

\subsection{Two Convergent Investigations}

The Studium Generale ORGANVM research programme has produced two
investigations that approach the formal structure of meaning from
complementary directions. RP-01, \emph{Toward a Grand Unified Semantics},
surveyed three traditions of formal semantics---logical (Tarski~\cite{Tarski1936};
Kripke~\cite{Kripke1959,Kripke1963}), programming language
(Scott and Strachey~\cite{ScottStrachey1971}; Hoare~\cite{Hoare1969}), and
natural language (Montague~\cite{Montague1973};
Heim and Kratzer~\cite{HeimKratzer1998})---and identified eight structural
bridges: compositionality, lambda calculus, the Curry-Howard-Lambek
correspondence, type theory, model theory, game semantics, dynamic semantics,
and proof-theoretic semantics.

RP-06, \emph{Grammars, Types, Proofs, and Categories}~\cite{Padavano2026RP06},
established a fourfold correspondence: at each level of the Chomsky hierarchy,
grammars correspond to automata (the classical result), but they also
correspond to regimes of type-theoretic expressiveness, and the convergence
reveals that grammars, types, proofs, and categories are structurally
isomorphic.

\subsection{The Question of Unification}

If meaning is a map from syntax to semantics, and if syntax has a unified
categorical structure, then meaning itself must have a unified categorical
structure---the structure of a functor from the syntactic category to the
semantic category.

\subsection{What Kind of Unification?}

\textbf{Kitcher's account}~\cite{Kitcher1981}: unification reduces the number
of independent explanatory patterns. The functorial pattern replaces multiple
independent descriptions of compositionality with a single argument pattern.
This is genuine but incremental.

\textbf{Morrison's taxonomy}~\cite{Morrison2000}: distinguishes theoretical
unification (Maxwell), reductive unification (Weinberg-Salam), and structural
unification. The present work achieves \emph{structural} unification: shared
mathematical architecture, not shared cause or substrate.

\textbf{The Erlangen Programme}~\cite{Klein1872}: Klein ``unified'' the
geometries of the nineteenth century by characterising each as the study of
invariants under a specific transformation group. As
Eilenberg and Mac Lane~\cite{EilenbergMacLane1945} noted, the categorical
perspective ``may be regarded as a continuation of the Klein Erlanger
Programm.''

\begin{itemize}
  \item \textbf{Claimed:} The three traditions share a common categorical
    architecture---the functorial passage from a syntactic category to a
    semantic category.
  \item \textbf{Not claimed:} Scientific unification in the Maxwell sense;
    novel empirical predictions; capture of pragmatics, embodied meaning, or
    non-compositional phenomena.
\end{itemize}

\subsection{The Central Claim}

\emph{The compositional, type-driven core of formal meaning---across the
logical, computational, and linguistic traditions---admits structural
unification under a single categorical architecture: the functorial passage
from a syntactic category to a semantic category.}

More precisely: there exists a family of categories $\mathbf{Syn}$ (syntactic
categories) and a \emph{fibred family} of categories $\mathbf{Sem}$ (semantic
categories, parametrised by interpretive tradition) such that meaning
assignment is a functor $F \colon \mathbf{Syn} \to \mathbf{Sem}$.

\subsection{Scope and Limitations}

The following fall outside the architecture's scope: pragmatics
(Grice~\cite{Grice1975}; Searle~\cite{Searle1969}); embodied meaning
(Lakoff and Johnson~\cite{LakoffJohnson1980};
Varela et al.~\cite{VarelaThompsonRosch1991}); non-compositional phenomena
(idioms, metaphor); dynamic semantics in its full generality
(Kamp~\cite{Kamp1981}; Heim~\cite{Heim1982};
Groenendijk and Stokhof~\cite{GroenendijkStokhof1991});
game-theoretic semantics; and probabilistic meaning.

\section{Recapitulation: The Eight Bridges and the Fourfold}

\subsection{The Bridge Taxonomy (RP-01)}

\begin{enumerate}
  \item \textbf{Compositionality} (Frege's Principle). Structural strength:
    \emph{isomorphism}.
  \item \textbf{Lambda calculus}. Strength: \emph{isomorphism}.
  \item \textbf{Curry-Howard-Lambek correspondence}. Strength:
    \emph{isomorphism} (three-way); \emph{homomorphism} (extension to natural
    language).
  \item \textbf{Type theory}. Strength: \emph{homomorphism}.
  \item \textbf{Model theory}. Strength: \emph{isomorphism} (classical
    fragments).
  \item \textbf{Game semantics}. Strength: \emph{homomorphism}.
  \item \textbf{Dynamic semantics}. Strength: \emph{analogy} tending toward
    homomorphism.
  \item \textbf{Proof-theoretic semantics}. Strength: \emph{homomorphism}.
\end{enumerate}

\subsection{The Fourfold Correspondence (RP-06)}

\begin{table}[ht]
\centering
\caption{The fourfold correspondence from RP-06}
\label{tab:fourfold}
\begin{tabular}{@{}llll@{}}
\toprule
Grammar & Logic & Computation & Category \\
\midrule
Syntactic category & Proposition & Type & Object \\
Derivation & Proof & Programme & Morphism \\
$A/B$ & Implication & Function type & Exponential \\
Application & Modus ponens & $\beta$-reduction & Evaluation \\
Lexicon & Axioms & Constants & Generators \\
\bottomrule
\end{tabular}
\end{table}

\subsection{Overlap and Divergence}

RP-01 and RP-06 converge on the Curry-Howard-Lambek correspondence, categorial
grammar, lambda calculus, and compositionality. They diverge on dynamic/game
semantics (RP-01 only), automata theory (RP-06 only), and the treatment of
distributional semantics.

\section{The Categorical Architecture}

\subsection{Meaning as Functor: The Fibred Perspective}

A functor $F \colon \mathbf{C} \to \mathbf{D}$ sends objects and morphisms of
$\mathbf{C}$ to $\mathbf{D}$, preserving composition ($F(g \circ f) =
F(g) \circ F(f)$) and identity ($F(\mathrm{id}_A) = \mathrm{id}_{F(A)}$).

The multiplicity of semantic categories---$\mathbf{Set}$ (classical logic),
$\mathbf{Dom}$ (Scott domains), $\mathbf{FdVect}$ (distributional
semantics), Kripke presheaves (modal logic)---means that the ``unified
functorial passage'' is a \emph{family} of functors parametrised by the
choice of semantic category.

\textbf{The fibred formulation.} Let $\mathbf{T}$ be the category of
\emph{interpretive traditions}. The semantic categories form a fibration
$p \colon \mathbf{Sem} \to \mathbf{T}$, where the fibre over each tradition
$t \in \mathbf{T}$ is the semantic category $\mathbf{Sem}_t$. The meaning
functor is a \emph{section} of this fibration: a family of functors
$F_t \colon \mathbf{Syn} \to \mathbf{Sem}_t$, coherent with respect to the
morphisms in $\mathbf{T}$~\cite{Grothendieck1971,Jacobs1999}.

\textbf{Why ``almost tautological'' is a feature.} That the functorial
perspective arises independently in logical semantics, programming language
semantics, and natural language semantics is evidence that the concept of
structure-preserving meaning assignment was already latent in all three
traditions. The Erlangen Programme was similarly ``tautological''---Klein's
contribution was to make the schema explicit.

\subsection{Compositionality as Functoriality}

The identification of compositionality with functoriality has been observed
by Janssen~\cite{Janssen1997}, Hodges~\cite{Hodges2001}, and others.
Montague's principle of compositionality asserts that there exists a
homomorphism $h \colon \mathbf{A}_{\mathrm{syn}} \to
\mathbf{A}_{\mathrm{sem}}$. This is precisely the assertion that meaning
assignment is a functor.

A failure of compositionality is a failure of the meaning map to be a
functor: a case in which $F(g \circ f) \neq F(g) \circ F(f)$. Idioms
provide the canonical example. Moggi's monadic
framework~\cite{Moggi1991} restores compositionality for effects by enriching
the semantic category.

\subsection{The Lambek Calculus as the Syntactic Category}

The Lambek calculus~\cite{Lambek1958} provides the canonical syntactic
category. As a category, it forms the \emph{free residuated monoid} generated
by the basic syntactic types. The residuation laws encode the adjunction:
\[
A \cdot B \to C \quad\text{iff}\quad A \to C / B \quad\text{iff}\quad B \to A \backslash C
\]

Its non-commutativity reflects word order. Its substructural character
connects it to linear logic~\cite{Girard1987}.

\subsection{Montague's Rule-to-Rule Correspondence as Natural Transformation}

Montague's rule-to-rule hypothesis can be formulated as a natural
transformation $\eta \colon S \Rightarrow M$ between syntactic and semantic
interpretation functors. The commutativity of the naturality square:

\begin{verbatim}
S(A) --eta_A--> M(A)
 |                |
S(f)             M(f)
 |                |
 v                v
S(B) --eta_B--> M(B)
\end{verbatim}

\noindent is precisely the rule-to-rule hypothesis: meaning assignment is
uniform and systematic. The intensional dimension adds further structure via
functors $\mathbf{Syn} \to [\mathbf{W}, \mathbf{Sem}]$, where $\mathbf{W}$ is
the category of possible worlds.

\subsection{The DisCoCat Framework: Promise, Proof of Concept, and Empirical Limits}

DisCoCat~\cite{CoeckeSadrzadehClark2010} defines meaning as a strong monoidal
functor $F \colon \mathbf{G} \to \mathbf{FdVect}$, where $\mathbf{G}$ is the
free rigid category generated by a pregroup grammar. This assigns vector
spaces to basic types and tensors to words.

\textbf{The empirical record: an honest assessment.} DisCoCat has been tested
on small-scale sentence similarity tasks. It has \emph{not} been demonstrated
on machine translation, open-domain QA, summarisation, dialogue, or any task
requiring reasoning over paragraphs. DisCoCat demonstrates that
reconciliation is feasible \emph{in principle}, for sentences of two or three
words. It is a proof of \emph{categorical concept}, not a practical bridge.

\textbf{Limitations.} (1)~$\mathbf{FdVect}$ cannot directly represent logical
operations (negation, quantification). (2)~Pregroup grammars are weakly
equivalent to context-free grammars. (3)~Tensor representations grow
exponentially with type arity.

\subsection{Compact Closed Categories and the Symmetric/Non-Symmetric Distinction}

Pregroup grammars form \emph{rigid} categories (non-symmetric compact closed),
while $\mathbf{FdVect}$ is \emph{symmetric} compact closed. The DisCoCat
functor maps from a non-symmetric structure to a symmetric one and therefore
\emph{cannot be fully faithful}: it collapses the distinction between left and
right adjoints. The structural preservation is weaker than a naive reading
would suggest, compensated by richer structure at the morphism level.

The compact closed structure provides a diagrammatic calculus (string
diagrams), imported from categorical quantum
mechanics~\cite{AbramskyCoecke2004}, which has led to quantum natural language
processing (QNLP) via the \texttt{lambeq} framework~\cite{Kartsaklis2021}.

\subsection{The Topos-Theoretic Perspective: Presheaf Semantics}

A presheaf over a category of contexts $\mathbf{C}$ is a functor
$F \colon \mathbf{C}^{\mathrm{op}} \to \mathbf{Set}$ that assigns to each
context a set of ``elements in context'' with coherent translations.

\textbf{Worked example: anaphora resolution as a sheaf condition.} Let the
category of contexts $\mathbf{C}$ be the category of discourse states (sets
of available discourse referents, ordered by inclusion). Define a presheaf
$R \colon \mathbf{C}^{\mathrm{op}} \to \mathbf{Set}$ assigning to each
discourse state the set of possible referent assignments.

For the discourse:
\begin{quote}
(1) ``A farmer owns a donkey.'' \\
(2) ``He beats it.''
\end{quote}

After sentence~(1), the discourse state is $D_1 = \{\text{farmer},
\text{donkey}\}$. The \emph{sheaf condition} says: if partial assignments
(resolving ``he'' and ``it'' separately) are compatible, there is a unique
total assignment extending both. When the sheaf condition \emph{fails}, we
get the formal diagnosis of an anaphoric failure. The ``donkey sentences'' of
dynamic semantics~\cite{Kamp1981,Heim1982} can be analysed in this framework.

\textbf{Connections to the architecture.} A topos is a cartesian closed
category (compositional semantics available within); the subobject classifier
provides non-classical truth values (vagueness, graded truth); the internal
logic is intuitionistic (connecting to Curry-Howard-Lambek).

\section{The Unification Claim}

\subsection{What Is Unified (and What Is Not)}

The architecture structurally unifies the compositional, type-driven core of
grammar, logic, computation, and natural language semantics. It does
\emph{not} unify pragmatic meaning, distributional meaning with
truth-conditional meaning (though DisCoCat shows the bridge is categorically
possible), formal structure with neural mechanisms, or compositional with
embodied approaches.

\subsection{The Five-Way Correspondence}

\begin{table}[ht]
\centering
\caption{The five-way correspondence}
\label{tab:fiveway}
\small
\begin{tabular}{@{}lllll@{}}
\toprule
Grammar & Logic & Computation & Category & Semantics \\
\midrule
Syntactic type & Proposition & Type & Object & Semantic type \\
Derivation & Proof & Programme & Morphism & Meaning-comp. \\
Concatenation & Conjunction & Product type & Product & Conjunction \\
Slash $A/B$ & $A \Rightarrow B$ & $A \to B$ & Exponential & Functional meaning \\
Sentence deriv. & Theorem proof & Well-typed term & Composed morph. & Sentence meaning \\
Lexicon & Axioms & Constants & Generators & Lexical meanings \\
Grammaticality & Provability & Well-typedness & Morph.\ existence & Meaningfulness \\
Ambiguity & Multiple proofs & Multiple terms & Multiple morph. & Multiple readings \\
Compositionality & Cut elimination & $\beta$-reduction & Composition law & Frege's principle \\
\bottomrule
\end{tabular}
\end{table}

\textbf{The correspondence as a commutative diagram.}

\begin{verbatim}
                          Syn
                        / | \ \
                       /  |  \ \
                      v   v   v  v
                   F_L  F_D  F_V  F_P
                    |    |    |    |
                    v    v    v    v
                 Sem_L Sem_D Sem_V Sem_P
                    \    |    |   /
                     \   |    |  /
                  U_LD \ U_DV | / U_VP
                       v v   v v
                 (natural transformations
                  witnessing coherence)
\end{verbatim}

\noindent where $\mathbf{Syn}$ is the syntactic category;
$F_L$, $F_D$, $F_V$, $F_P$ are the logical, domain-theoretic, distributional,
and proof-theoretic meaning functors; and the natural transformations $U$ at
the base are the formal expression of the bridges identified by RP-01.

\subsection{Evidence For}

\textbf{Formal:} Curry-Howard, Lambek's theorem, Pentus's equivalence,
Montague's rule-to-rule, adequacy/full abstraction results.
\textbf{Structural:} The same mathematical concepts appear with the same
structural character across traditions. \textbf{Historical:} The pattern of
convergent discovery (Chomsky/Backus on CFGs; Curry/Howard on
proofs-as-programmes; game semantics independently in three fields).

\subsection{Evidence Against}

Pragmatics and non-compositionality; neural representations and opacity;
concurrency and interaction (Curry-Howard for concurrency remains open);
the gap between formal adequacy and empirical coverage.

\section{The Distributional Challenge}

\subsection{Transformers: Confronting Success Without Categorical Structure}

Transformers demonstrate that meaning can be computed (or closely approximated)
without implementing the categorical architecture. This is not a minor
observation. The architecture offers structural transparency, guaranteed
compositionality, and a framework for comparison. Transformers offer empirical
performance, scalability, and learning from data.

Probing studies~\cite{HewittManning2019,Clark2019,Jawahar2019} show that
transformers encode syntactic dependency relations. ``Partially compatible'' is
a weak claim: the representations are distributed and continuous, not the
discrete algebraic structures of the Lambek calculus.

\textbf{The honest framing:} the architecture and transformers may describe the
same phenomena at different levels of abstraction, or different phenomena that
co-occur in language. The question is empirical and not yet decisive.

\subsection{Testable Predictions}

\textbf{Prediction 1:} Transformers should show more systematic,
structure-sensitive behaviour on compositional than non-compositional
constructions.

\textbf{Prediction 2:} Architectures encoding categorical structure should
outperform standard transformers on compositional generalisation tasks
(COGS~\cite{KimLinzen2020}, SCAN~\cite{LakeBaroni2018}).

\textbf{Prediction 3:} Compositional patterns in one tradition should have
precise analogues in others, predictable via the five-way correspondence.

\textbf{Prediction 4:} Words with more complex grammatical types should
require higher-rank tensors in DisCoCat, reflected in learning difficulty.

\section{Discussion}

\subsection{What the Architecture Achieves}

Structural transparency (a single framework connecting RP-01 and RP-06);
precision (informal ``analogies'' replaced by functors, natural
transformations, adjunctions); a framework for systematic comparison of
semantic traditions via the fibration.

\subsection{What the Architecture Does Not Achieve}

It is not scientific unification. It does not compete with transformers. It
does not capture pragmatics, embodied meaning, non-compositional phenomena,
or game-theoretic meaning.

\subsection{The Value of Structural Unification}

The Erlangen Programme is universally regarded as a major intellectual
achievement, despite being ``merely'' structural. The analogous claim: the
functorial architecture organises the study of formal meaning, clarifies
relationships between traditions, and enables systematic transfer
(e.g., monadic techniques from PL semantics for non-compositionality in NL).
The evidence of DisCoCat, of the \texttt{lambeq} framework, and of the
testable predictions suggests the architecture is productive.

\section{Conclusion}

This paper has synthesised two prior investigations into a single categorical
architecture of compositional formal meaning. The architecture has seven
components: meaning as functor (fibred), compositionality as functoriality,
the Lambek calculus as syntactic category, Montague's rule-to-rule as natural
transformation, DisCoCat (honestly assessed), compact closed categories (with
the symmetric/non-symmetric distinction), and presheaf semantics (with a
worked example of anaphora resolution).

The central claim---the five-way correspondence among derivation, proof,
programme, morphism, and meaning-composition---has been stated precisely and
evaluated against supporting and countervailing evidence. The evidence
supports the claim for the classical, compositional, typed core.

This is structural unification in the tradition of Klein's Erlangen Programme.
The three traditions of compositional formal semantics are three perspectives
on a single mathematical structure. The functorial architecture makes this
structure explicit. Whether making structure explicit constitutes
``unification'' depends on one's philosophy of unification. That it
constitutes a productive intellectual achievement---in the tradition of the
Erlangen Programme---the present author submits with confidence.

\section*{Acknowledgments}

This work was conducted as part of the Studium Generale ORGANVM research
program. The author acknowledges the use of AI language models (Claude,
Anthropic) as research assistants for literature synthesis, structural
drafting, and adversarial review via the Triadic Review Protocol. All
intellectual claims, arguments, and editorial decisions are the author's own.

\medskip
\noindent\textcopyright~2026 A.~Padavano. This work is licensed under a
Creative Commons Attribution 4.0 International License (CC BY 4.0).

\begin{thebibliography}{99}

\bibitem{AbramskyCoecke2004}
S.~Abramsky and B.~Coecke.
\newblock A categorical semantics of quantum protocols.
\newblock In \emph{Proc.\ 19th IEEE Symp.\ Logic in Computer Science}, pp.~415--425, 2004.

\bibitem{Ajdukiewicz1935}
K.~Ajdukiewicz.
\newblock Die syntaktische Konnexitat.
\newblock \emph{Studia Philosophica}, 1:1--27, 1935.

\bibitem{BarHillel1953}
Y.~Bar-Hillel.
\newblock A quasi-arithmetical notation for syntactic description.
\newblock \emph{Language}, 29(1):47--58, 1953.

\bibitem{Carpenter1997}
B.~Carpenter.
\newblock \emph{Type-Logical Semantics}.
\newblock MIT Press, 1997.

\bibitem{Chomsky1956}
N.~Chomsky.
\newblock Three models for the description of language.
\newblock \emph{IRE Trans.\ Information Theory}, 2(3):113--124, 1956.

\bibitem{Church1940}
A.~Church.
\newblock A formulation of the simple theory of types.
\newblock \emph{J.\ Symbolic Logic}, 5(2):56--68, 1940.

\bibitem{Clark2019}
K.~Clark, U.~Khandelwal, O.~Levy, and C.~D.~Manning.
\newblock What does BERT look at?
\newblock In \emph{Proc.\ 2019 ACL Workshop BlackboxNLP}, pp.~276--286, 2019.

\bibitem{CoeckeSadrzadehClark2010}
B.~Coecke, M.~Sadrzadeh, and S.~Clark.
\newblock Mathematical foundations for a compositional distributional model of meaning.
\newblock \emph{Linguistic Analysis}, 36(1--4):345--384, 2010.

\bibitem{CoquandHuet1988}
T.~Coquand and G.~Huet.
\newblock The Calculus of Constructions.
\newblock \emph{Information and Computation}, 76(2--3):95--120, 1988.

\bibitem{CurryFeys1958}
H.~B.~Curry and R.~Feys.
\newblock \emph{Combinatory Logic}, Vol.~I.
\newblock North-Holland, 1958.

\bibitem{Davidson1967}
D.~Davidson.
\newblock Truth and meaning.
\newblock \emph{Synthese}, 17(1):304--323, 1967.

\bibitem{EilenbergMacLane1945}
S.~Eilenberg and S.~Mac Lane.
\newblock General theory of natural equivalences.
\newblock \emph{Trans.\ Amer.\ Math.\ Soc.}, 58(2):231--294, 1945.

\bibitem{Girard1987}
J.-Y.~Girard.
\newblock Linear logic.
\newblock \emph{Theoretical Computer Science}, 50(1):1--101, 1987.

\bibitem{Grice1975}
H.~P.~Grice.
\newblock Logic and conversation.
\newblock In P.~Cole and J.~L.~Morgan, eds., \emph{Syntax and Semantics}, Vol.~3, pp.~41--58. Academic Press, 1975.

\bibitem{GroenendijkStokhof1991}
J.~Groenendijk and M.~Stokhof.
\newblock Dynamic predicate logic.
\newblock \emph{Linguistics and Philosophy}, 14(1):39--100, 1991.

\bibitem{Grothendieck1971}
A.~Grothendieck.
\newblock \emph{Rev\^{e}tements \'{e}tales et groupe fondamental} (SGA 1).
\newblock LNMS~224, Springer, 1971.

\bibitem{Heim1982}
I.~Heim.
\newblock \emph{The Semantics of Definite and Indefinite Noun Phrases}.
\newblock PhD dissertation, University of Massachusetts, 1982.

\bibitem{HeimKratzer1998}
I.~Heim and A.~Kratzer.
\newblock \emph{Semantics in Generative Grammar}.
\newblock Blackwell, 1998.

\bibitem{HewittManning2019}
J.~Hewitt and C.~D.~Manning.
\newblock A structural probe for finding syntax in word representations.
\newblock In \emph{Proc.\ NAACL-HLT 2019}, pp.~4129--4138, 2019.

\bibitem{Hoare1969}
C.~A.~R.~Hoare.
\newblock An axiomatic basis for computer programming.
\newblock \emph{Communications of the ACM}, 12(10):576--580, 1969.

\bibitem{Hodges2001}
W.~Hodges.
\newblock Formal features of compositionality.
\newblock \emph{J.\ Logic, Language and Information}, 10(1):7--28, 2001.

\bibitem{Howard1980}
W.~A.~Howard.
\newblock The formulae-as-types notion of construction.
\newblock In J.~P.~Seldin and J.~R.~Hindley, eds., \emph{To H.~B.~Curry}, pp.~479--490. Academic Press, 1980. (Manuscript 1969.)

\bibitem{Jacobs1999}
B.~Jacobs.
\newblock \emph{Categorical Logic and Type Theory}.
\newblock Elsevier, 1999.

\bibitem{Janssen1997}
T.~M.~V.~Janssen.
\newblock Compositionality.
\newblock In J.~van Benthem and A.~ter Meulen, eds., \emph{Handbook of Logic and Language}, pp.~417--473. Elsevier, 1997.

\bibitem{Jawahar2019}
G.~Jawahar, B.~Sagot, and D.~Seddah.
\newblock What does BERT learn about the structure of language?
\newblock In \emph{Proc.\ 57th ACL}, pp.~3651--3657, 2019.

\bibitem{Kamp1981}
H.~Kamp.
\newblock A theory of truth and semantic representation.
\newblock In J.~Groenendijk et al., eds., \emph{Formal Methods in the Study of Language}, pp.~277--322. Mathematisch Centrum, 1981.

\bibitem{Kartsaklis2021}
D.~Kartsaklis, I.~Fan, R.~Yeung, A.~Pearson, R.~Lorenz, A.~Toumi, G.~de~Felice, K.~Meichanetzidis, S.~Clark, and B.~Coecke.
\newblock lambeq: An efficient high-level Python library for quantum NLP.
\newblock arXiv:2110.04236, 2021.

\bibitem{KimLinzen2020}
N.~Kim and T.~Linzen.
\newblock COGS: A compositional generalization challenge based on semantic interpretation.
\newblock In \emph{Proc.\ EMNLP 2020}, pp.~9087--9105, 2020.

\bibitem{Kitcher1981}
P.~Kitcher.
\newblock Explanatory unification.
\newblock \emph{Philosophy of Science}, 48(4):507--531, 1981.

\bibitem{Klein1872}
F.~Klein.
\newblock Vergleichende Betrachtungen \"{u}ber neuere geometrische Forschungen.
\newblock \emph{Mathematische Annalen}, 43:63--100, 1893. (Original 1872.)

\bibitem{Kripke1959}
S.~Kripke.
\newblock A completeness theorem in modal logic.
\newblock \emph{J.\ Symbolic Logic}, 24(1):1--14, 1959.

\bibitem{Kripke1963}
S.~Kripke.
\newblock Semantical considerations on modal logic.
\newblock \emph{Acta Philosophica Fennica}, 16:83--94, 1963.

\bibitem{LakeBaroni2018}
B.~M.~Lake and M.~Baroni.
\newblock Generalization without systematicity.
\newblock In \emph{Proc.\ ICML 2018}, pp.~2873--2882, 2018.

\bibitem{LakoffJohnson1980}
G.~Lakoff and M.~Johnson.
\newblock \emph{Metaphors We Live By}.
\newblock Univ.\ Chicago Press, 1980.

\bibitem{Lambek1958}
J.~Lambek.
\newblock The mathematics of sentence structure.
\newblock \emph{American Mathematical Monthly}, 65(3):154--170, 1958.

\bibitem{Lambek1968}
J.~Lambek.
\newblock Deductive systems and categories I.
\newblock \emph{Mathematical Systems Theory}, 2:287--318, 1968.

\bibitem{Lambek1999}
J.~Lambek.
\newblock Type grammar revisited.
\newblock In A.~Lecomte et al., eds., \emph{Logical Aspects of Computational Linguistics}, LNCS~1582, pp.~1--27. Springer, 1999.

\bibitem{MartinLof1984}
P.~Martin-L\"{o}f.
\newblock \emph{Intuitionistic Type Theory}.
\newblock Bibliopolis, 1984.

\bibitem{Moggi1991}
E.~Moggi.
\newblock Notions of computation and monads.
\newblock \emph{Information and Computation}, 93(1):55--92, 1991.

\bibitem{Montague1973}
R.~Montague.
\newblock The proper treatment of quantification in ordinary English.
\newblock In K.~J.~J.~Hintikka et al., eds., \emph{Approaches to Natural Language}, pp.~221--242. Reidel, 1973.

\bibitem{Moortgat1997}
M.~Moortgat.
\newblock Categorial type logics.
\newblock In J.~van Benthem and A.~ter Meulen, eds., \emph{Handbook of Logic and Language}, pp.~93--177. Elsevier, 1997.

\bibitem{Morrison2000}
M.~Morrison.
\newblock \emph{Unifying Scientific Theories}.
\newblock Cambridge Univ.\ Press, 2000.

\bibitem{Padavano2026RP06}
A.~Padavano.
\newblock The fourfold correspondence: Grammar, automaton, type, and proof from Chomsky to Curry-Howard-Lambek.
\newblock arXiv preprint, 2026. (Forthcoming.)

\bibitem{ScottStrachey1971}
D.~S.~Scott and C.~Strachey.
\newblock Toward a mathematical semantics for computer languages.
\newblock In J.~Fox, ed., \emph{Proc.\ Symp.\ Computers and Automata}, pp.~19--46. Polytechnic Inst.\ Brooklyn Press, 1971.

\bibitem{Searle1969}
J.~R.~Searle.
\newblock \emph{Speech Acts}.
\newblock Cambridge Univ.\ Press, 1969.

\bibitem{Steedman2000}
M.~Steedman.
\newblock \emph{The Syntactic Process}.
\newblock MIT Press, 2000.

\bibitem{Tarski1936}
A.~Tarski.
\newblock The concept of truth in formalized languages.
\newblock In A.~Tarski, \emph{Logic, Semantics, Metamathematics}, 2nd ed.\ (1983), pp.~152--278. Hackett.

\bibitem{VarelaThompsonRosch1991}
F.~J.~Varela, E.~Thompson, and E.~Rosch.
\newblock \emph{The Embodied Mind}.
\newblock MIT Press, 1991.

\end{thebibliography}

\end{document}
